\documentclass{ximera}

\input{../../../preamble.tex}


\definecolor{fvdnavy}{HTML}{071D43}
\definecolor{fvdcyan}{HTML}{20BFC6}
\definecolor{fvdcyanlight}{HTML}{EFFCFB}
\definecolor{fvdmagenta}{HTML}{F10D69}
\definecolor{fvdmagentalight}{HTML}{FFF1F7}
\definecolor{fvdtext}{HTML}{163044}





%=================================================
% Pedagogische omgevingen
% PDF: tcolorbox
% HTML: learning-box
%=================================================

\newenvironment{voorkennis}
{%
    \ifdefined\HCode
        \HCode{%
<section class="learning-box learning-box--prior">
<div class="learning-box__header">
<span class="learning-box__icon" aria-hidden="true"></span>
<span class="learning-box__title">Voorkennis</span>
</div>
<div class="learning-box__content">
        }%
        \begin{itemize}
    \else
       \begin{tcolorbox}[
    enhanced,
    breakable,
    colback=fvdcyanlight,
    colframe=fvdcyan,
    coltext=fvdtext,
    boxrule=0.5pt,
    leftrule=4pt,
    arc=4mm,
    outer arc=4mm,
    left=7mm,
    right=6mm,
    top=5mm,
    bottom=5mm,
    before skip=6mm,
    after skip=6mm,
    title=Voorkennis,
    coltitle=fvdcyan,
    fonttitle=\bfseries\Large,
    attach title to upper={\par\smallskip},
    boxed title style={
        empty,
        left=0mm,
        right=0mm,
        top=0mm,
        bottom=0mm
    },
    overlay={
        \draw[
            fvdcyan,
            line width=0.5pt
        ]
        ([xshift=42mm,yshift=-13mm]frame.north west)
        --
        ([xshift=-7mm,yshift=-13mm]frame.north east);
    }
]
        \begin{itemize}
    \fi
}
{%
    \end{itemize}
    \ifdefined\HCode
        \HCode{%
</div>
</section>
        }%
    \else
        \end{tcolorbox}
    \fi
}


\newenvironment{leerdoelen}
{%
    \ifdefined\HCode
        \HCode{%
<section class="learning-box learning-box--goals">
<div class="learning-box__header">
<span class="learning-box__icon" aria-hidden="true"></span>
<span class="learning-box__title">Leerdoelen</span>
</div>
<div class="learning-box__content">
        }%
        \begin{itemize}
    \else
\begin{tcolorbox}[
    enhanced,
    breakable,
    colback=fvdmagentalight,
    colframe=fvdmagenta,
    coltext=fvdtext,
    boxrule=0.5pt,
    leftrule=4pt,
    arc=4mm,
    outer arc=4mm,
    left=7mm,
    right=6mm,
    top=5mm,
    bottom=5mm,
    before skip=6mm,
    after skip=6mm,
    title=Leerdoelen,
    coltitle=fvdmagenta,
    fonttitle=\bfseries\Large,
    attach title to upper={\par\smallskip},
    boxed title style={
        empty,
        left=0mm,
        right=0mm,
        top=0mm,
        bottom=0mm
    },
    overlay={
        \draw[
            fvdmagenta,
            line width=0.5pt
        ]
        ([xshift=42mm,yshift=-13mm]frame.north west)
        --
        ([xshift=-7mm,yshift=-13mm]frame.north east);
    }
]
        \begin{itemize}
    \fi
}
{%
    \end{itemize}
    \ifdefined\HCode
        \HCode{%
</div>
</section>
        }%
    \else
        \end{tcolorbox}
    \fi
}


\addPrintStyle{../../..}

\title{Oefeningen — Producten en quotiënten afleiden}
\author{Ingmar Herreman}

\begin{abstract}
In deze oefeningen oefen je de productregel en de quotiëntregel.

Je leert niet alleen formules toepassen, maar ook herkennen welke
strategie het meest geschikt is, resultaten vereenvoudigen en fouten
in redeneringen opsporen.

De oefeningen met  vormen samen het minimumtraject.
\end{abstract}

\begin{document}

% FVD_CHAPTER_DATA_START
% Automatisch gegenereerd uit index.tex.
% Niet handmatig aanpassen.
\fvdchapterdata{3}{Afleiden van samengestelde functie}{1}
% FVD_CHAPTER_DATA_END










\maketitle




\FVDbeginExercises


% ------------------------------------------------------------
% OEFENING 1
% ------------------------------------------------------------

\begin{oefening}[Productregel rechtstreeks toepassen]{\oefB\ESS}

Bepaal telkens de afgeleide.

\begin{question}
\[
f(x)=(x^2+1)(x-3).
\]

\begin{oplossing}
Met
\[
u=x^2+1,
\qquad
v=x-3
\]
geldt
\[
u'=2x,
\qquad
v'=1.
\]

Dus
\[
\begin{aligned}
f'(x)
&=2x(x-3)+(x^2+1)\\
&=3x^2-6x+1.
\end{aligned}
\]

\[
\boxed{f'(x)=3x^2-6x+1}
\]
\end{oplossing}
\end{question}


\begin{question}
\[
g(x)=(2x^3-x)(x^2+4).
\]

\begin{oplossing}
\[
\begin{aligned}
g'(x)
&=(6x^2-1)(x^2+4)
 +(2x^3-x)(2x)\\
&=10x^4+21x^2-4.
\end{aligned}
\]

\[
\boxed{g'(x)=10x^4+21x^2-4}
\]
\end{oplossing}
\end{question}


\begin{question}
\[
h(x)=x^2(3x-5).
\]

Bereken de afgeleide met de productregel.

\begin{oplossing}
\[
\begin{aligned}
h'(x)
&=2x(3x-5)+3x^2\\
&=9x^2-10x.
\end{aligned}
\]

\[
\boxed{h'(x)=9x^2-10x}
\]
\end{oplossing}
\end{question}

\end{oefening}


% ------------------------------------------------------------
% OEFENING 2
% ------------------------------------------------------------

\begin{oefening}[Product of eerst uitwerken?]{\oefB\ESS}

Bepaal telkens de afgeleide op de meest efficiënte manier.

\begin{question}
\[
f(x)=x(x^2+4x-1).
\]

\begin{oplossing}
Eerst uitwerken is eenvoudig:
\[
f(x)=x^3+4x^2-x.
\]

Dus
\[
\boxed{f'(x)=3x^2+8x-1}.
\]
\end{oplossing}
\end{question}


\begin{question}
\[
g(x)=(x+2)(x-5).
\]

\begin{oplossing}
\[
g(x)=x^2-3x-10,
\]
dus
\[
\boxed{g'(x)=2x-3}.
\]
\end{oplossing}
\end{question}


\begin{question}
\[
h(x)=(x^2+1)(x^3-2).
\]

\begin{oplossing}
Hier is de productregel handig:
\[
\begin{aligned}
h'(x)
&=2x(x^3-2)+(x^2+1)3x^2\\
&=5x^4+3x^2-4x.
\end{aligned}
\]
\end{oplossing}
\end{question}

\end{oefening}


% ------------------------------------------------------------
% OEFENING 3
% ------------------------------------------------------------

\begin{oefening}[Quotiëntregel]{\oefB\ESS}

Bepaal telkens eerst het domein en daarna de afgeleide.

\begin{question}
\[
f(x)=\frac{x^2+3}{x-1}.
\]

\begin{oplossing}
\[
\operatorname{dom}f=\mathbb{R}\setminus\{1\}.
\]

\[
\begin{aligned}
f'(x)
&=
\frac{2x(x-1)-(x^2+3)}{(x-1)^2}\\
&=
\frac{x^2-2x-3}{(x-1)^2}\\
&=
\frac{(x-3)(x+1)}{(x-1)^2}.
\end{aligned}
\]
\end{oplossing}
\end{question}


\begin{question}
\[
g(x)=\frac{2x-1}{x^2+1}.
\]

\begin{oplossing}
Omdat
\[
x^2+1>0,
\]
is
\[
\operatorname{dom}g=\mathbb{R}.
\]

\[
g'(x)
=
\frac{-2x^2+2x+2}{(x^2+1)^2}.
\]
\end{oplossing}
\end{question}


\begin{question}
\[
h(x)=\frac{x^3+x}{x^2-4}.
\]

\begin{oplossing}
\[
\operatorname{dom}h
=
\mathbb{R}\setminus\{-2,2\}.
\]

\[
h'(x)
=
\frac{x^4-13x^2-4}{(x^2-4)^2}.
\]
\end{oplossing}
\end{question}

\end{oefening}


% ------------------------------------------------------------
% OEFENING 4
% ------------------------------------------------------------

\begin{oefening}[Let op de volgorde]{\oefB\ESS}

Een leerling gebruikt voor
\[
f(x)=\frac{x^2+1}{x-3}
\]
de formule

\[
f'(x)=
\frac{
(x^2+1)\cdot1-2x(x-3)
}{
(x-3)^2
}.
\]

\begin{question}
Welke fout heeft de leerling gemaakt?

\begin{oplossing}
De volgorde van de teller van de quotiëntregel werd omgekeerd.

De correcte teller is
\[
u'v-uv'.
\]
\end{oplossing}
\end{question}

\begin{question}
Bereken de correcte afgeleide.

\begin{oplossing}
\[
\begin{aligned}
f'(x)
&=
\frac{
2x(x-3)-(x^2+1)
}{
(x-3)^2
}\\
&=
\boxed{
\frac{x^2-6x-1}{(x-3)^2}
}.
\end{aligned}
\]
\end{oplossing}
\end{question}

\end{oefening}


% ------------------------------------------------------------
% OEFENING 5
% ------------------------------------------------------------

\begin{oefening}[Niet elke breuk vraagt om de quotiëntregel]{\oefB\ESS}

Herschrijf eerst en bepaal daarna de afgeleide.

\begin{question}
\[
f(x)=\frac{4}{x^3}.
\]

\begin{oplossing}
\[
f(x)=4x^{-3},
\]
dus
\[
\boxed{
f'(x)=-12x^{-4}=-\frac{12}{x^4}
}.
\]
\end{oplossing}
\end{question}


\begin{question}
\[
g(x)=\frac{5x^3}{x}.
\]

\begin{oplossing}
Voor $x\neq0$:
\[
g(x)=5x^2.
\]

Dus
\[
\boxed{g'(x)=10x\qquad(x\neq0)}.
\]

De oorspronkelijke domeinvoorwaarde $x\neq0$ blijft gelden.
\end{oplossing}
\end{question}


\begin{question}
\[
h(x)=\frac{x^4-2x^2}{x^2}.
\]

\begin{oplossing}
Voor $x\neq0$:
\[
h(x)=x^2-2.
\]

Dus
\[
\boxed{h'(x)=2x\qquad(x\neq0)}.
\]
\end{oplossing}
\end{question}

\end{oefening}


% ------------------------------------------------------------
% OEFENING 6
% ------------------------------------------------------------

\begin{oefening}[Productregel combineren met gekende regels]{\oefB}

Bepaal telkens de afgeleide en vereenvoudig.

\begin{question}
\[
f(x)=3x^2(x^3-2x+1).
\]

\begin{oplossing}
\[
\begin{aligned}
f'(x)
&=6x(x^3-2x+1)
 +3x^2(3x^2-2)\\
&=15x^4-18x^2+6x.
\end{aligned}
\]
\end{oplossing}
\end{question}


\begin{question}
\[
g(x)=(x^2-4)(2x^2+x+1).
\]

\begin{oplossing}
\[
\begin{aligned}
g'(x)
&=2x(2x^2+x+1)
 +(x^2-4)(4x+1).
\end{aligned}
\]

Een vereenvoudigde vorm is
\[
\boxed{
g'(x)=8x^3+3x^2-14x-4
}.
\]
\end{oplossing}
\end{question}

\end{oefening}


% ------------------------------------------------------------
% OEFENING 7
% ------------------------------------------------------------

\begin{oefening}[Quotiënt vereenvoudigen]{\oefB\ESS}

Bepaal de afgeleide en schrijf je antwoord zo eenvoudig mogelijk.

\begin{question}
\[
f(x)=\frac{x^2}{x+1}.
\]

\begin{oplossing}
\[
\begin{aligned}
f'(x)
&=\frac{2x(x+1)-x^2}{(x+1)^2}\\
&=\boxed{\frac{x(x+2)}{(x+1)^2}}.
\end{aligned}
\]
\end{oplossing}
\end{question}


\begin{question}
\[
g(x)=\frac{x^2-1}{x}.
\]

\begin{oplossing}
Herschrijven is eenvoudiger:
\[
g(x)=x-\frac1x.
\]

Dus
\[
\boxed{
g'(x)=1+\frac1{x^2}
}
\qquad(x\neq0).
\]
\end{oplossing}
\end{question}

\end{oefening}


% ============================================================
% VERDIEPING
% ============================================================

\FVDendExercises



\FVDbeginExercises


% ------------------------------------------------------------
% OEFENING 8
% ------------------------------------------------------------

\begin{oefening}[Welke strategie kies je?]{\oefU\ESS}

Bepaal de afgeleide.

Schrijf bij elke functie eerst welke strategie je kiest:
\begin{itemize}
  \item eerst vereenvoudigen;
  \item productregel;
  \item quotiëntregel.
\end{itemize}

Motiveer je keuze kort.

\begin{question}
\[
f(x)=\frac{6}{x^4}.
\]
\end{question}

\begin{question}
\[
g(x)=(x^2+3)(2x^3-1).
\]
\end{question}

\begin{question}
\[
h(x)=\frac{x^2+4}{2x-1}.
\]
\end{question}

\begin{question}
\[
p(x)=\frac{x^5+x^3}{x^3}.
\]
\end{question}

\begin{oplossing}
Mogelijke keuzes:

\[
f(x)=6x^{-4}
\]
dus herschrijven.

Voor $g$ is de productregel natuurlijk.

Voor $h$ is de quotiëntregel natuurlijk.

Voor $p$ geldt, met $x\neq0$,
\[
p(x)=x^2+1,
\]
dus eerst vereenvoudigen.
\end{oplossing}

\end{oefening}


% ------------------------------------------------------------
% OEFENING 9
% ------------------------------------------------------------

\begin{oefening}[Twee methoden vergelijken]{\oefU\ESS}

Gegeven:
\[
f(x)=\frac{x^3-2x^2}{x}.
\]

\begin{question}
Bepaal $f'$ met de quotiëntregel.

\begin{oplossing}
\[
f'(x)
=
\frac{
(3x^2-4x)x-(x^3-2x^2)
}{
x^2
}.
\]

Vereenvoudigen geeft
\[
f'(x)=2x-2
\qquad (x\neq0).
\]
\end{oplossing}
\end{question}


\begin{question}
Vereenvoudig eerst $f$ en bepaal daarna $f'$.

\begin{oplossing}
Voor $x\neq0$:
\[
f(x)=x^2-2x.
\]

Dus
\[
f'(x)=2x-2.
\]
\end{oplossing}
\end{question}


\begin{question}
Welke methode verkies je hier? Motiveer.

\begin{oplossing}
Eerst vereenvoudigen is hier korter en minder foutgevoelig.
De domeinvoorwaarde $x\neq0$ moet wel behouden blijven.
\end{oplossing}
\end{question}

\end{oefening}


% ------------------------------------------------------------
% OEFENING 10
% ------------------------------------------------------------

\begin{oefening}[Fouten analyseren]{\oefU\ESS}

Een leerling schrijft:

\[
\begin{aligned}
f(x)&=(x^2+1)(x^3-2),\\
f'(x)&=2x(3x^2).
\end{aligned}
\]

\begin{question}
Leg uit waarom dit fout is.

\begin{oplossing}
De leerling heeft de afgeleiden van beide factoren met elkaar
vermenigvuldigd.

Maar
\[
(uv)'=u'v+uv'.
\]
\end{oplossing}
\end{question}


\begin{question}
Bereken de correcte afgeleide.

\begin{oplossing}
\[
\begin{aligned}
f'(x)
&=2x(x^3-2)+(x^2+1)3x^2\\
&=5x^4+3x^2-4x.
\end{aligned}
\]
\end{oplossing}
\end{question}

\end{oefening}


% ------------------------------------------------------------
% OEFENING 11
% ------------------------------------------------------------

\begin{oefening}[Een afgeleide controleren zonder opnieuw te beginnen]{\oefU}

Een leerling beweert dat voor
\[
f(x)=\frac{x^2+1}{x+1}
\]
geldt:
\[
f'(x)=\frac{x^2+2x+1}{(x+1)^2}.
\]

\begin{question}
Factoriseer de teller van het voorgestelde antwoord.

\begin{oplossing}
\[
x^2+2x+1=(x+1)^2.
\]

Het voorgestelde antwoord zou dus gelijk zijn aan $1$.
\end{oplossing}
\end{question}


\begin{question}
Kan $f'(x)$ identiek gelijk zijn aan $1$?
Controleer door $f$ eerst via veeltermdeling te herschrijven.

\begin{oplossing}
\[
\frac{x^2+1}{x+1}
=
x-1+\frac{2}{x+1}.
\]

Dus
\[
f'(x)
=
1-\frac{2}{(x+1)^2}.
\]

Het voorgestelde antwoord is dus fout.
\end{oplossing}
\end{question}

\end{oefening}


% ------------------------------------------------------------
% OEFENING 12
% ------------------------------------------------------------

\begin{oefening}[Van afgeleide naar verloop]{\oefU\ESS}

Gegeven:
\[
f(x)=\frac{x^2}{x+1}.
\]

Je weet dat
\[
f'(x)=\frac{x(x+2)}{(x+1)^2}.
\]

\begin{question}
Welke waarden mogen niet in het domein voorkomen?

\begin{oplossing}
\[
x\neq-1.
\]
\end{oplossing}
\end{question}


\begin{question}
Welke factor bepaalt het teken van de noemer van $f'$?

\begin{oplossing}
De noemer
\[
(x+1)^2
\]
is positief voor alle $x\neq-1$.
\end{oplossing}
\end{question}


\begin{question}
Waar is $f'(x)=0$?

\begin{oplossing}
\[
x(x+2)=0
\]
dus
\[
x=-2
\qquad\text{of}\qquad
x=0.
\]
\end{oplossing}
\end{question}


\begin{question}
Stel een tekentabel voor $f'$ op en bepaal waar $f$ stijgt
en daalt.

\begin{oplossing}
Het teken wordt bepaald door
\[
x(x+2).
\]

Dus:
\[
f'>0
\quad\text{voor}\quad
x<-2
\quad\text{en}\quad
x>0,
\]
en
\[
f'<0
\quad\text{voor}\quad
-2<x<-1
\quad\text{en}\quad
-1<x<0.
\]

Daarom stijgt $f$ op
\[
]-\infty,-2[
\quad\text{en}\quad
]0,+\infty[
\]
en daalt ze op
\[
]-2,-1[
\quad\text{en}\quad
]-1,0[.
\]
\end{oplossing}
\end{question}

\end{oefening}


% ------------------------------------------------------------
% OEFENING 13
% ------------------------------------------------------------

\begin{oefening}[Parameter en horizontale raaklijn]{\oefU\ESS}

Beschouw
\[
f_a(x)=\frac{x^2+a}{x-1}.
\]

\begin{question}
Bepaal $f_a'(x)$.

\begin{oplossing}
\[
f_a'(x)
=
\frac{x^2-2x-a}{(x-1)^2}.
\]
\end{oplossing}
\end{question}


\begin{question}
Bepaal $a$ zodat de grafiek in $x=2$ een horizontale raaklijn heeft.

\begin{oplossing}
Een horizontale raaklijn vereist
\[
f_a'(2)=0.
\]

Dus
\[
4-4-a=0,
\]
waaruit
\[
\boxed{a=0}.
\]
\end{oplossing}
\end{question}

\end{oefening}


% ------------------------------------------------------------
% OEFENING 14
% ------------------------------------------------------------

\begin{oefening}[Parallel geschakelde weerstanden]{\oefU}

Een vaste weerstand $R>0$ wordt parallel geschakeld met een
variabele weerstand $x>0$.

De vervangingsweerstand is
\[
R_{\mathrm{v}}(x)
=
\frac{Rx}{R+x}.
\]

\begin{question}
Bepaal
\[
R_{\mathrm{v}}'(x).
\]

\begin{oplossing}
\[
\begin{aligned}
R_{\mathrm{v}}'(x)
&=
\frac{
R(R+x)-Rx
}{
(R+x)^2
}\\
&=
\boxed{
\frac{R^2}{(R+x)^2}
}.
\end{aligned}
\]
\end{oplossing}
\end{question}


\begin{question}
Welk teken heeft $R_{\mathrm{v}}'(x)$ voor $R>0$ en $x>0$?

\begin{oplossing}
\[
R_{\mathrm{v}}'(x)>0.
\]
\end{oplossing}
\end{question}


\begin{question}
Wat betekent dit fysisch?

\begin{oplossing}
Wanneer de variabele weerstand $x$ groter wordt, neemt ook de
vervangingsweerstand toe.

De afgeleide geeft bovendien aan hoe gevoelig de vervangingsweerstand
is voor een kleine verandering van $x$.
\end{oplossing}
\end{question}

\end{oefening}


% ============================================================
% GEVORDERD
% ============================================================

\FVDendExercises



\FVDbeginExercises


% ------------------------------------------------------------
% OEFENING 15
% ------------------------------------------------------------

\begin{oefening}[Reconstrueer een parameter]{\oefG}

Gegeven:
\[
f_a(x)=\frac{x^2+a}{x+2}.
\]

De raaklijn aan de grafiek in het punt met abscis $x=0$
heeft richtingscoëfficiënt $-1$.

Bepaal $a$.

\begin{oplossing}
\[
f_a'(x)
=
\frac{
2x(x+2)-(x^2+a)
}{
(x+2)^2
}.
\]

Dus
\[
f_a'(0)
=
-\frac{a}{4}.
\]

Uit
\[
-\frac{a}{4}=-1
\]
volgt
\[
\boxed{a=4}.
\]
\end{oplossing}

\end{oefening}


% ------------------------------------------------------------
% OEFENING 16
% ------------------------------------------------------------

\begin{oefening}[Ontwerp zelf een functie]{\oefG}

Zoek een rationale functie van de vorm
\[
f(x)=\frac{x^2+a}{x+b}
\]
waarvoor:

\[
x=1
\]
niet tot het domein behoort en de grafiek in $x=0$ een horizontale
raaklijn heeft.

\begin{question}
Bepaal eerst $b$.

\begin{oplossing}
Omdat $x=1$ uit het domein moet verdwijnen:
\[
1+b=0,
\]
dus
\[
b=-1.
\]
\end{oplossing}
\end{question}


\begin{question}
Bepaal daarna $a$.

\begin{oplossing}
Nu:
\[
f(x)=\frac{x^2+a}{x-1}.
\]

De afgeleide is
\[
f'(x)
=
\frac{x^2-2x-a}{(x-1)^2}.
\]

Voor een horizontale raaklijn in $x=0$:
\[
f'(0)=0.
\]

Dus
\[
-a=0,
\]
waaruit
\[
a=0.
\]

Een mogelijke functie is dus
\[
\boxed{
f(x)=\frac{x^2}{x-1}
}.
\]
\end{oplossing}
\end{question}

\end{oefening}


% ------------------------------------------------------------
% OEFENING 17
% ------------------------------------------------------------

\begin{oefening}[Een algemene eigenschap onderzoeken]{\oefG}

Beschouw
\[
f(x)=\frac{ax+b}{cx+d},
\]
waarbij
\[
cx+d\neq0.
\]

\begin{question}
Bepaal $f'(x)$.

\begin{oplossing}
\[
\begin{aligned}
f'(x)
&=
\frac{
a(cx+d)-c(ax+b)
}{
(cx+d)^2
}\\
&=
\frac{
acx+ad-acx-bc
}{
(cx+d)^2
}\\
&=
\boxed{
\frac{ad-bc}{(cx+d)^2}
}.
\end{aligned}
\]
\end{oplossing}
\end{question}


\begin{question}
Welke bijzondere eigenschap heeft de teller van $f'(x)$?

\begin{oplossing}
De teller
\[
ad-bc
\]
is constant en hangt niet af van $x$.
\end{oplossing}
\end{question}


\begin{question}
Wat kun je zeggen over het teken van $f'$ als
\[
ad-bc>0?
\]

\begin{oplossing}
Omdat
\[
(cx+d)^2>0
\]
op het domein, geldt overal
\[
f'(x)>0.
\]

De functie is dus op elk interval van haar domein strikt stijgend.
\end{oplossing}
\end{question}


\begin{question}
Wat gebeurt er als
\[
ad-bc=0?
\]

\begin{oplossing}
Dan is
\[
f'(x)=0
\]
op elk interval van het domein.

In dat geval zijn teller en noemer evenredig en stelt de breuk,
waar gedefinieerd, een constante functie voor.
\end{oplossing}
\end{question}

\end{oefening}


% ------------------------------------------------------------
% OEFENING 18
% ------------------------------------------------------------

\begin{oefening}[Bewering beoordelen]{\oefG}

Een leerling beweert:

\begin{quote}
Een functie die als quotiënt geschreven is, moet altijd met de
quotiëntregel worden afgeleid.
\end{quote}

Onderzoek deze bewering.

\begin{question}
Geef een tegenvoorbeeld waarvoor herschrijven duidelijk eenvoudiger is.

\begin{oplossing}
Bijvoorbeeld:
\[
f(x)=\frac{3}{x^2}=3x^{-2}.
\]

Dan is
\[
f'(x)=-6x^{-3}.
\]
\end{oplossing}
\end{question}


\begin{question}
Geef een voorbeeld waarvoor de quotiëntregel wel een natuurlijke
keuze is.

\begin{oplossing}
Bijvoorbeeld:
\[
g(x)=\frac{x^2+1}{x-2}.
\]
\end{oplossing}
\end{question}


\begin{question}
Formuleer een betere algemene strategie.

\begin{oplossing}
Onderzoek eerst of de functie eenvoudig kan worden herschreven of
vereenvoudigd. Kies pas daarna de afleidingsregel die tot de meest
efficiënte en overzichtelijke berekening leidt.
\end{oplossing}
\end{question}

\end{oefening}


\FVDendExercises




\end{document}